\documentclass[12pt,a4paper]{article}
\usepackage{cmap} % Makes the PDF copiable. See http://tex.stackexchange.com/a/64198/25761
\usepackage[T1]{fontenc}
\usepackage[brazil]{babel}
\usepackage[utf8]{inputenc}
\usepackage{amsmath}
\usepackage{amsfonts}
\usepackage{amssymb}
\usepackage{amsthm}
\usepackage{textcomp} % \degree
\usepackage{gensymb} % \degree
\usepackage[usenames,svgnames,dvipsnames]{xcolor}
\usepackage{hyperref}
\usepackage{graphicx}
\usepackage[margin=2cm]{geometry}

\hypersetup{
    colorlinks = true,
    allcolors = {blue}
}

% TODO: Consider using exsheets
% http://linorg.usp.br/CTAN/macros/latex/contrib/exsheets/exsheets_en.pdf
%
% http://ctan.org/tex-archive/macros/latex/contrib/exercise/
% Options: answerdelayed,lastexercise,noanswer
\usepackage[answerdelayed,lastexercise]{exercise}

\addto\captionsbrazil{%
\def\listexercisename{Lista de exerc\'icios}%
\def\ExerciseName{Exerc\'icio}%
\def\AnswerName{Solu\c{c}\~ao do exerc\'icio}%
\def\ExerciseListName{Ex.}%
\def\AnswerListName{Solu\c{c}\~ao}%
\def\ExePartName{Parte}%
\def\ArticleOf{de\ }%
}

\renewcommand{\ExerciseHeaderTitle}{(\ExerciseTitle)\ }
\renewcommand{\ExerciseListHeader}{%\ExerciseHeaderDifficulty%
\textbf{%\ExerciseListName\
\ExerciseHeaderNB.\ %
%\ --- \
\ExerciseHeaderTitle}%
%\ExerciseHeaderOrigin
\ignorespaces}
\renewcommand{\AnswerListHeader}{\textbf{\ExerciseHeaderNB.\ (\AnswerListName)\ }}

\newcommand{\fixme}{{\color{red}(...)}}
\newcommand*\diff{\mathop{}\!\mathrm{d}}

\renewcommand{\theenumi}{\alph{enumi}}
\renewcommand\labelenumi{(\theenumi) }

\newcommand*\tipo{Prova II}
\newcommand*\turma{TADS121-01C}
\newcommand*\disciplina{CDI0001}
\newcommand*\eu{Helder G. G. de Lima}
\newcommand*\data{07/10/2016}

\author{\eu}
\title{\tipo - \disciplina}
\date{\data}

\begin{document}
\thispagestyle{empty}
\newgeometry{margin=2cm,bottom=0.5cm}
\begin{center}
\includegraphics[width=9.0cm]{marca} \\
\textbf{\tipo\ (\disciplina / \turma)} \\
Prof. \eu\footnote{
Este é um material de acesso livre distribuído sob os termos da licença \href{https://creativecommons.org/licenses/by-sa/4.0/deed.pt_BR}{Creative Commons BY-SA 4.0}.}
\end{center}

\noindent Nome do(a) aluno(a): \underline{\hspace{9,7cm}} Data: \underline{\data}

%\section*{Instruções}
\begin{center}\fbox{
\begin{minipage}{14cm}

{\footnotesize
\begin{itemize}
\renewcommand{\theenumi}{\Roman{enumi}}
\item Identifique-se em todas as folhas.
\item Mantenha o celular e os demais equipamentos eletrônicos desligados durante a prova.
\item Escolha as questões que irá resolver de modo a totalizar até 10 pontos.
\end{itemize}
}

\end{minipage}
}
\end{center}


%\section*{Questões}
\begin{ExerciseList}
\Exercise%[title={2,0}]
Considere a função $f$ definida por $y = f(x) = \sqrt{4x + 1}$.
\begin{enumerate}
\item \textbf{(1,0)} Calcule $f^\prime(2)$, por meio do limite que define esta derivada.
\item \textbf{(0,5)} Utilize as regras de derivação para obter $f^\prime(a) = \dfrac{ \diff{y} } { {dx} } (a)$, e verifique que para $a = 2$, o valor obtido coincide com o resultado do item anterior.
\item \textbf{(0,5)} Qual é o domínio da função $f^\prime$, isto é, quais são todos os pontos em que $f$ é derivável?
\item \textbf{(2,0)} Utilize a técnica de derivação implícita para calcular $\dfrac{ \diff{x} } { {dy} }$ no ponto $y$ (isto é, obtenha a derivada da função $x=f^{-1}(y)$ definida implicitamente por $y = \sqrt{4x + 1}$).
\end{enumerate}

\Answer \fixme

\Exercise[title={2,0}] Seja $g$ a função definida por $g(x) = \ln(12x-3x^2)$. Obtenha a equação da reta tangente ao gráfico de $g$ no ponto $P=(3, g(3))$.
\Answer \fixme

\Exercise%[title={2,0}]
Seja $h(x)=(e^x + 1)^2$.
\begin{enumerate}
\item \textbf{(1,0)} Calcule $h^\prime(x)$, $h^{\prime\prime}(x)$ e $h^{(3)}(x)$.
\item \textbf{(0,5)} Escreva a fórmula geral para a $n$-ésima derivada de $h(x)$.
\item \textbf{(0,5)} Utilize a fórmula do item anterior para calcular $h^{(7)}(0)$.
\end{enumerate}
\Answer \fixme

\Exercise[title={2,0}] Utilize uma aproximação por diferenciais (aproximação linear) para estimar o volume de material necessário para fazer uma bola com espessura de $0,1 cm$ e raio interno de $10 cm$.

\Exercise[title={2,0}] Ao andar por uma calçada, você está se afastando de uma luminária de $2,7m$ de altura a uma velocidade de $1,5 m/s$. Assumindo que você tenha $1,8m$ de altura, utilize o que aprendeu sobre derivadas para determinar com que velocidade o comprimento da sua sombra no chão aumenta conforme você se afasta da luminária.

\end{ExerciseList}

\begin{center}
BOA PROVA!
\end{center}

%\newpage
%\restoregeometry
%\section*{Respostas}
%\shipoutAnswer
\end{document}
